\section{Distance geometry in matrix space}
\label{sec:geometry}

The row-sum identity $\M\mathbf 1=\Net$ gives knowledge-matrix perturbations an
exact visible/invisible accounting. This accounting is shared with any
representation carrying the same fixed row-sum constraint --- in particular
per-class gradient$\times$input (Theorem~\ref{thm:weff}); what it excludes is
\emph{stored activation-based} representations, for which no fixed
parameter-independent vector plays the role of $\mathbf 1$
(Proposition~\ref{prop:dichotomy}).

\begin{theorem}[Visible/invisible decomposition]\label{thm:pyth}
For any $\Delta\M\in\R^{C\times(d+1)}$ with logit displacement
$\Delta\Net=\Delta\M\,\mathbf 1$, let $P=\Delta\Net\,\mathbf 1^{\!\top}/(d{+}1)$
and $Q=\Delta\M-P$. Then $P$ is the orthogonal projection of $\Delta\M$ onto
constant-row matrices and the unique minimum-norm matrix with row sums
$\Delta\Net$; $Q\mathbf 1=0$; and
\[
\|\Delta\M\|_F^2=\frac{\|\Delta\Net\|_2^2}{d+1}+\|Q\|_F^2,
\qquad\text{hence}\qquad
\|\Delta\M\|_F\;\ge\;\frac{\|\Delta\Net\|_2}{\sqrt{d+1}},
\]
with equality iff $\Delta\M$ has constant rows. We call
$\rho_{\mathrm{vis}}=\|\Delta\Net\|^2/((d{+}1)\|\Delta\M\|_F^2)$ the visible
fraction and $A:=(d{+}1)\,\rho_{\mathrm{vis}}=(d_\Psi/d_M)^2$ the
\textbf{coherence} of the perturbation; both are computable from the two
stored scalars $(d_M,d_\Psi)$, and rank statistics of $d_M/d_\Psi$ are exactly
reverse rank statistics of $A$.
\end{theorem}

\begin{clarifybox}
\emph{``Rank statistics of $d_M/d_\Psi$ are exactly reverse rank statistics of
$A$.''} Since $A=(d_\Psi/d_M)^2=1/r^2$ with $r=d_M/d_\Psi>0$, and $t\mapsto1/t^2$ is
strictly decreasing on $(0,\infty)$, the map from one to the other is
order-reversing pair by pair. Any statistic that sees the sample only through
the ordering of its values --- order statistics, quantiles, medians, Spearman
and Kendall coefficients, the concordance $W$ --- therefore transfers between
$r$ and $A$ with its direction flipped and nothing else changed. In
particular the attack-family ordering we report is one and the same finding
whichever of the two we tabulate, which is exactly why that ordering does not
borrow anything from $A$ having a validated meaning. (The single caveat is
arithmetic rather than conceptual: an interpolated median of an even-count
sample averages two order statistics, so it commutes with the transform only
up to the gap between them --- below our quoted precision; see
Appendix~\ref{app:proofs}.)
\end{clarifybox}

\begin{clarifybox}
\emph{Why $A=(d_\Psi/d_M)^2$, and not the ratio $d_\Psi/d_M$ itself?} Four reasons,
of which only the first is essential. \emph{(a)} Theorem~\ref{thm:pyth} is an
\emph{energy} statement: $\|\Delta\M\|_F^2$ splits \emph{additively} into
$\|\Delta\Net\|^2/(d{+}1)$ and $\|Q\|_F^2$ because the two components are
orthogonal, and additivity across orthogonal components is a property of
squared norms, not of norms. The dimensionless quantity the theorem actually
hands us is therefore a share of squared norm, namely $\rho_{\mathrm{vis}}$;
$d_\Psi/d_M$ is not a share of anything. \emph{(b)} $A$ \emph{is} that share,
renormalised: $A=(d{+}1)\rho_{\mathrm{vis}}$, the factor $(d{+}1)$ chosen so
that the single-pixel case reads $A=1$ instead of $1/(d{+}1)$. \emph{(c)} In
the squared scale the two theorem-given reference lines are the clean numbers
$A=1$ and $A\le d$; in the unsquared scale they would be $1$ and $\sqrt d$,
and the additive split would not be legible at all. \emph{(d)} Nothing is
lost: the two are strictly monotone functions of each other, so every
rank-based conclusion is identical under either, as the previous note records.
\end{clarifybox}

$Q$ is invisible \emph{to the evaluated logit displacement}: it encodes real
changes of the local linearisation that the endpoint logits cannot see. The
split is an exact orthogonal decomposition --- short to prove but load-bearing,
since the entire descriptive geometry below rests on it. The
decomposition is invariant under output rescaling, under the gauge group, and
under per-coordinate input reparameterisations (pixel-unit changes move
nothing); it is \emph{not} invariant under general input rotations --- the
statistic is tied to the pixel basis, which for images is the natural one, and
we make no claim of rotation invariance.
This theorem retires our earlier amplification claims: raw Frobenius
comparisons inflate with mass-spreading and per-coordinate (RMS) comparisons
deflate with dimension, and both are monotone transforms of the same
unit-free coherence $A=(d_\Psi/d_M)^2$ (a monotone transform of the $d_M/d_\Psi$
ratio). We use $A$ as a \emph{geometric descriptor} of a perturbation, read
against the two theorem-given reference lines $A=1$ (the one-pixel law,
Theorem~\ref{thm:within}(ii)) and $A\le d$ (the within-region cap,
Theorem~\ref{thm:within}(i)). We do \emph{not} tie $A$ to detectability or to
any decision-relevant outcome: it is descriptive geometry, not a validated
statistic, and the standalone empirical finding it is reported alongside is the
attack-family rank ordering, which does not depend on $A$ having any such link.

\begin{theorem}[Within-region anatomy]\label{thm:within}
If $x$ and $y=x+\delta$ share the strict mask pattern, then
$\M(y)-\M(x)=[\,J\,\mathrm{diag}(\delta)\mid 0\,]$ --- the bias column cancels
exactly --- and $d_M^2=\sum_i\delta_i^2\|J_{:,i}\|_2^2$. Consequently:
\emph{(i)} (coherence cap) $A\le d$, with the sharp maximum iff all weighted
columns $\delta_iJ_{:,i}$ are equal and nonzero; \emph{(ii)} (one-pixel law)
for $\delta=s\,e_{i_0}$ with $J_{:,i_0}\neq0$: $d_M=d_\Psi$ exactly, i.e.\
$A=1$, independent of pixel and magnitude; \emph{(iii)} ($k$-sparse bound) a
perturbation supported on $k$ pixels has at most $k$ nonzero columns,
$A\le k$, and column-energy participation ratio at most $k$.
\end{theorem}

\begin{clarifybox}
\emph{Column-energy participation ratio.} Write $e_i=\delta_i^2\|J_{:,i}\|_2^2$
for the energy that column $i$ contributes to $d_M^2=\sum_ie_i$. The
participation ratio of that energy profile is
$\mathrm{PR}=\bigl(\sum_ie_i\bigr)^2\big/\sum_ie_i^2$ --- the standard
inverse-participation-ratio measure of how many columns actually carry the
energy. It equals $k$ exactly when $k$ columns share the energy equally and
the rest are silent, and $1$ when a single column carries all of it, so it
reads as an \emph{effective number of active columns}. Part (iii) then says a
$k$-sparse perturbation cannot have more than $k$ of them, which is immediate:
only the $k$ columns in the support have $e_i\neq0$.
\end{clarifybox}

\begin{clarifybox}
\emph{The one-pixel law, and why $A=1$ does not mean ``nothing happened''.}
The algebra is two lines. Within a region
$\M(y)-\M(x)=[\,J\,\mathrm{diag}(\delta)\mid0\,]$, so
$d_M^2=\sum_i\delta_i^2\|J_{:,i}\|^2$; for $\delta=s\,e_{i_0}$ this is
$d_M=|s|\,\|J_{:,i_0}\|$, while $\Delta\Net=J\delta=s\,J_{:,i_0}$ gives
$d_\Psi=|s|\,\|J_{:,i_0}\|$ as well. Hence $d_M=d_\Psi$ and $A=1$, for every pixel
and every magnitude.

What is pinned is the \emph{ratio}, not the effect. Both $d_M$ and $d_\Psi$ grow
linearly in $|s|$, so a large single-pixel edit moves the matrix and the
logits a great deal; it is their quotient that cannot move. $A$ measures
interference among the weighted columns $\delta_iJ_{:,i}$, and a one-pixel
displacement switches on exactly one of them --- there is nothing to reinforce
and nothing to cancel, so all of the matrix motion is logit-visible by
default. $A=1$ is the no-interference baseline, not a claim that the
perturbation is inert.

The pixel-by-pixel worry then dissolves in two steps. \emph{First}, $A$ is a
functional of the displacement $\delta=y-x$, not of a path. Editing $k$ pixels
one at a time produces $k$ single-pixel displacements each with $A=1$, but the
object this theorem scores is the \emph{cumulative} $\delta$, which is
$k$-sparse and not one-sparse; $A$ is not additive along a path and does not
accumulate. For the cumulative $\delta$, part (iii) allows any value up to
$k$: above $1$ when the $\delta_iJ_{:,i}$ reinforce, below $1$ when they
cancel, and near $1$ when they are mutually incoherent. \emph{Second}, a long
walk leaves the region. Once $x$ and $y$ no longer share a mask pattern the
hypothesis here fails and Theorem~\ref{thm:anatomy} governs instead: the
crossing dyads satisfy $K\mathbf 1=0$, so they add energy to $d_M$ while
contributing exactly nothing to $d_\Psi$, and generically push $A$ down. This is
the regime the experiments live in --- full-image adversarial displacements,
crossing many walls, are measured at median $A\le0.23$
(Section~\ref{sec:study2}).
\end{clarifybox}

\begin{theorem}[Smooth $+$ crossing anatomy]\label{thm:anatomy}
For a generic pair $(x,y)$ whose segment crosses region walls transversally at
$t_1<\dots<t_N$, with $\bar J=\int_0^1J(x+t\delta)\,dt$,
\[
\M(y)-\M(x)=\bigl[\,\bar J\,\mathrm{diag}(\delta)\mid 0\,\bigr]
\;+\;K,\qquad
K=\sum_{j}\sigma_j\,u_{k_j}\bigl[\,v_{k_j}^{\top}\mathrm{diag}(z_j)\mid
\gamma_{k_j}\,\bigr],\qquad K\mathbf 1=0,
\]
where each crossing contributes a rank-one dyad whose row sums vanish (the
flipped unit's pre-activation is zero at the crossing), and
$\Delta\Net=\bar J\,\delta$: \emph{the crossing part moves the germ and not the
function}. The bias column of $\M(y)-\M(x)$ equals $c(y)-c(x)
=\sum_j\sigma_j\gamma_{k_j}u_{k_j}$ exactly; a nonzero bias column therefore
certifies at least one crossing (no false positives; the converse fails ---
bias-free networks are silent). The endpoint mask-Hamming $H$ lower-bounds the
crossing count, with parity equality, at one forward pass per endpoint.
\end{theorem}

\begin{clarifybox}
\emph{The symbols in $K$, and what a rank-one dyad is.} A \emph{dyad} is an
outer product $uv^{\top}$ of two vectors: a matrix of rank one, every column a
multiple of $u$ and every row a multiple of $v^{\top}$. Each wall crossing
contributes exactly one such term because a single unit flipping changes the
region's Jacobian by a rank-one update --- the unit has one way in from the
input and one way out to the logits. Fix the $j$-th crossing, at which unit
$k_j$ of some hidden layer $\ell^{*}$ flips (dyad lemma,
Appendix~\ref{app:proofs}). Then $z_j=x+t_j\delta$ is the point on the segment
at which that crossing occurs, so $\mathrm{diag}(z_j)$ is the input scaling
that turns germ data into knowledge-matrix data \emph{there};
$\sigma_j\in\{+1,-1\}$ records the direction of the flip ($0\to1$ gives $+1$,
$1\to0$ gives $-1$), which for a ReLU unit is exactly the jump in its slope
across the wall; $u_{k_j}\in\R^{C}$ is the output-side vector, the column
through which unit $k_j$ reaches the $C$ logits along the downstream masked
product; and $v_{k_j}\in\R^{d}$, $\gamma_{k_j}\in\R$ are the input-side data,
defined by writing the unit's pre-activation as an affine function of the
network input, $z^{(\ell^{*})}_{k_j}(x')=v_{k_j}^{\top}x'+\gamma_{k_j}$, so
that $v_{k_j}^{\top}$ is a row of the upstream masked product and
$\gamma_{k_j}$ the bias accumulated below layer $\ell^{*}$. The wall is the
hyperplane $\{v_{k_j}^{\top}x'+\gamma_{k_j}=0\}$, which is precisely why the
dyad's row sums vanish: $[\,v^{\top}\mathrm{diag}(z_j)\mid\gamma\,]\mathbf 1
=v^{\top}z_j+\gamma=0$ at the crossing point.
\end{clarifybox}

\begin{clarifybox}
\emph{The path-averaged Jacobian $\bar J$ is a familiar object.}
$\bar J=\int_0^1J(x+t\delta)\,dt$ is the mean value of the Jacobian along the
segment from $x$ to $y$, and the identity $\Delta\Net=\bar J\delta$ is the mean
value theorem for vector-valued maps --- equivalently the fundamental theorem
of calculus along the segment --- which is exact here because $\Net$ is Lipschitz
and piecewise affine on it. It is also, exactly, the object behind
\emph{integrated gradients} \citep{sundararajan2017axiomatic}: the
integrated-gradients attribution of $y$ against baseline $x$ is
$(y-x)\odot\int_0^1\nabla\Net\bigl(x+t(y-x)\bigr)\,dt$, whose value for class
$c$ and coordinate $i$ is the $(c,i)$ entry of $\bar J\,\mathrm{diag}(\delta)$
--- the smooth block above. Their completeness axiom
$\sum_i\mathrm{IG}_i=\Nety-\Netx$ is the row-sum statement
$\bar J\delta=\Delta\Net$, and the axiomatic basis they invoke is the Aumann--Shapley
average-gradient cost-sharing rule. This sharpens rather than weakens the
equivalence we claim in Theorem~\ref{thm:weff}: at a single input the
knowledge matrix is per-class gradient$\times$input plus a bias column, and
\emph{between} two inputs its smooth part is integrated gradients. What the
decomposition adds is the remainder integrated gradients has no name for ---
the crossing term $K$, invisible to the endpoint logits and therefore
invisible to any completeness axiom phrased in terms of them.
\end{clarifybox}

\begin{clarifybox}
\emph{``The crossing part moves the germ and not the function'' does not
contradict Definition~\ref{def:germ}.} The germ is a \emph{local} object
attached to a \emph{single} point, and it is determined there by the function
(Theorem~\ref{thm:germ}); nothing in this theorem makes the germ at a fixed
point ambiguous. The sentence is about the \emph{displacement between two
different points}, and its ``function'' means the endpoint logit displacement
$\Delta\Net=\Nety-\Netx$ --- not the function as a whole. The claim is that
$\Delta\Net$ is accounted for in full by the smooth term, since
$\Delta\Net=\bar J\delta$ while $K\mathbf 1=0$, so the crossings contribute
nothing to it. They are nevertheless exactly what changes the local affine
data between the endpoints: $J(y)-J(x)=\sum_j\sigma_ju_{k_j}v_{k_j}^{\top}$
and $c(y)-c(x)=\sum_j\sigma_j\gamma_{k_j}u_{k_j}$, both carried by $K$ and
neither legible in $\Delta\Net$. The long-hand reading is therefore: between $x$
and $y$ the germ moves by more than the endpoint output values can reveal.
Both germs remain determined by the same single function, at two different
points --- which is Theorem~\ref{thm:germ}, not a violation of it.
\end{clarifybox}

\begin{clarifybox}
\emph{The endpoint mask-Hamming distance $H$, and ``one forward pass per
endpoint''.} $H$ counts the hidden units whose $0/1$ activation mask differs
between the two endpoints,
$H=\sum_{\ell,k}\bigl|D^{(\ell)}_k(x)-D^{(\ell)}_k(y)\bigr|$. It lower-bounds
the number of wall crossings along the segment, with equality of parity: a
unit that ends up flipped must have crossed an odd number of times, a unit
that ends up unflipped an even number --- possibly zero, possibly two, and a
unit that crosses and crosses back is invisible to the endpoints. Equality
$H=N$ holds iff no unit flips twice. The reason it is worth stating is cost:
$H$ needs only the two mask records that evaluating the network at $x$ and at
$y$ already produces --- one forward pass per endpoint, two in total --- with
no need to trace the segment, locate the crossing times $t_j$, or identify
which units flipped. It is the cheapest available certificate that two inputs
lie in different activation regions.
\end{clarifybox}

\begin{conjecture}[Mechanism of the attack-family ordering]\label{conj:mechanism}
Iterative small-step attacks cross fewer walls and align with high-energy
Jacobian columns, raising $A$; one-shot sign-based and gradient-free attacks
do the opposite. Testable as a size-controlled negative rank correlation
between per-pair $H$ and $A$. A pilot on per-pair VGG data (attack-success
filtered, size-controlled) finds the predicted sign on all three attacks
tested (Table~\ref{tab:mechanism-pilot}); the exact-$\|\delta\|$-controlled
replication on the full architecture set is future work.
\end{conjecture}

\begin{clarifybox}
\emph{``High-energy Jacobian columns.''} The energy of column $i$ of the
Jacobian is $\|J_{:,i}\|_2^2$: the weight with which a unit change in pixel
$i$ reaches the logits. A perturbation \emph{aligns with high-energy columns}
when its mass $|\delta_i|$ sits on the coordinates where $\|J_{:,i}\|$ is
large. The link to $A$ runs through Theorem~\ref{thm:within}: writing
$a_i=\delta_iJ_{:,i}$, we have $A=\|\sum_ia_i\|^2/\sum_i\|a_i\|^2$, so $A$
records how far the weighted columns reinforce rather than cancel --- maximal
($A=d$) when they are all equal, and $A=1$ whenever a single one stands alone,
which is the one-pixel law. Concentrating $\delta$ on a few high-energy
columns therefore suppresses the many small, mutually cancelling contributions
that pull $A$ below $1$, while by Theorem~\ref{thm:anatomy} each wall crossing
avoided removes a dyad that adds to $d_M$ and nothing to $d_\Psi$. The conjecture
is that both halves point the same way for iterative small-step attacks; it
remains a conjecture, and the pilot tests only its rank-correlation
consequence, not the mechanism itself.
\end{clarifybox}

\begin{proposition}[Penultimate distances carry no germ-level accounting]
\label{prop:dichotomy}
The decomposition of Theorem~\ref{thm:pyth} exists because
$\M\mathbf 1=\Net$ holds with a fixed, parameter-independent vector $\mathbf 1$
(it is shared by any representation with the same row-sum constraint, e.g.\
gradient$\times$input, Theorem~\ref{thm:weff}). Penultimate
features admit no analogue: on a gauge orbit the stored distance $d_h$ sweeps
an unbounded range while $\Net$, $d_\Psi$ and $d_M$ are fixed, so any gauge-invariant
statistic built from stored penultimate norms is independent of $d_h$ ---
i.e.\ not a statistic of the penultimate representation at all.
\end{proposition}

\begin{clarifybox}
\emph{Attribution: the fact underneath this proposition is not ours.} That a
positive per-neuron rescaling leaves a ReLU network's function exactly
unchanged while moving its hidden activations --- so that activation norms,
and any distance built from them, can be rescaled essentially at will --- is
the positive-homogeneity (rescaling) symmetry of ReLU networks. It is used as
an optimisation-geometry tool by \citet{neyshabur2015path} in Path-SGD and,
most familiarly, by \citet{dinh2017sharp} to show that the sharpness of a
minimum is reparameterisation-dependent; it is the same symmetry that appears
here as the gauge group, and as the rescaling half of
Lemma~\ref{lem:quiver-inv}. What this proposition contributes is only the
consequence for the accounting of Theorem~\ref{thm:pyth}: because no
gauge-invariant statistic can be built from stored penultimate norms, the
visible/invisible split has no penultimate analogue --- there is no fixed,
parameter-independent vector playing the role of $\mathbf 1$ on the feature
side. The symmetry is prior work; the corollary drawn from it for the
decomposition is what we claim.
\end{clarifybox}

\subsection{Teleportation: what is bounded and what must be measured}
\label{sec:teleport-claims}

\begin{theorem}[No $C^0$ bound]\label{thm:nogo}
For every $\varepsilon,B>0$ there exist same-architecture networks $\Net$ and
$g$ with $\sup_{x'}\|\Netxp-g(x')\|\le\varepsilon$, identical masks at $x$, and
$\|\M_g(x)-\M_\Psi(x)\|_F\ge B$. Function-space closeness does not control germ
drift.
\end{theorem}

\clarify{$C^0$ closeness means closeness in \emph{value}, uniformly: two maps are
$C^0$-close when $\sup_{x'}\|\Netxp-g(x')\|\le\varepsilon$, the metric of the
space $C^0$ of continuous functions under the supremum norm. The superscript
counts derivatives, so $C^1$ closeness would additionally require the
derivatives to agree. The theorem says the first buys nothing about the second,
and the reason is that the knowledge matrix records the germ --- first-order
data. The standard witness is $g=\Net+\varepsilon\sin(\cdot/\varepsilon^{2})$,
which stays uniformly within $\varepsilon$ of $\Net$ everywhere while its
derivative differs by $1/\varepsilon$. This is why we report the drift in
Section~\ref{sec:study1} as a measurement rather than inferring it from the
logits: agreement of the logits to any tolerance is, by itself, no evidence at
all about the matrices.}

\begin{proposition}[Three-tier claim structure for teleportation]\label{prop:visdrift}
\emph{(i)} Under an exact isomorphism, the knowledge-matrix drift is zero.
\emph{(ii)} Under an approximate teleportation with logit gate $\varepsilon$,
the \emph{visible} drift is bounded unconditionally:
$\|(\M_\tau(x)-\M(x))\mathbf 1\|=\|\Psi_\tau(W\!,f)(x)-\Netx\|\le\varepsilon$.
\emph{(iii)} The \emph{invisible} drift admits no a priori bound from the gate
(Theorem~\ref{thm:nogo}). A conditional bound exists when
both networks are affine on a common ball of radius $r$ around $x$:
$\|\M_\tau(x)-\M(x)\|_F\le(\varepsilon_r/r)\sqrt{\min(C,d)}\,\|x\|_\infty
+\varepsilon_r(1+\|x\|_2/r)$.
\end{proposition}

\begin{clarifybox}
\emph{The logit gate; gate-bounded versus gate-pinned.} The \emph{logit gate}
of a transform $\tau$ at $x$ is the measured discrepancy between the two
networks' outputs at the same input,
$\varepsilon=\|\Psi_\tau(W\!,f)(x)-\Netx\|$. It is an empirical quantity, not a
theoretical one: an experiment measures it and then uses it as an acceptance
threshold, treating $\tau$ as function-preserving when the gate falls below a
stated tolerance. A quantity is \emph{gate-bounded} when the theory bounds it
by $\varepsilon$, and \emph{gate-pinned} when the theory forces it to
\emph{equal} $\varepsilon$ as an algebraic identity. Part (ii) is the
gate-pinned case: the row-sum identity gives
$\|(\M_\tau(x)-\M(x))\mathbf 1\|=\|\Psi_\tau(W\!,f)(x)-\Netx\|=\varepsilon$ exactly, with
no inequality anywhere, so a table of visible drift is a table of the gate
itself and says nothing about invariance beyond what the gate already says.
Part (iii) is the genuinely unbounded case: by Theorem~\ref{thm:nogo} the gate
constrains the invisible drift not at all, which is why the conditional bound
there has to carry the extra hypothesis that both networks are affine on a
common ball. Whether a given transform actually needs this three-tier
treatment is settled empirically in Section~\ref{sec:study1b}.
\end{clarifybox}

Our teleportation study (Section~\ref{sec:study1b}) sits at tier~(i): the
transform is an exact isomorphism, batch normalisation in evaluation mode
included (Appendix~\ref{app:teleport-why}), so both components of the drift
vanish in exact arithmetic and what the experiment reports is the resolution of
its arithmetic. We measure both components rather than asserting either from a
theorem: the visible one through identity~(ii), which is an equality, and the
invisible one directly, at $0.8$--$2.5\times10^{-15}$ relative
(Appendix~\ref{app:teleport-km}). Tiers~(ii) and~(iii) are stated because they
are what one needs for a transform that genuinely is approximate; the no-go of
Theorem~\ref{thm:nogo} stands regardless, being a statement about what $C^0$
closeness cannot control rather than about any particular transform. The empirical
visible/invisible split is exhibited separately, on adversarial pairs where the
gate is exact, in Study~2 (Section~\ref{sec:study2}).
